\documentclass[]{article}


% Copyright (c) 2026 Simon Crase

% This script is free software: you can redistribute it and/or modify
% it under the terms of the GNU General Public License as published by
% the Free Software Foundation, either version 3 of the License, or
% (at your option) any later version.

% This script is distributed in the hope that it will be useful,
% but WITHOUT ANY WARRANTY; without even the implied warranty of
% MERCHANTABILITY or FITNESS FOR A PARTICULAR PURPOSE.  See the
% GNU General Public License for more details.

% You should have received a copy of the GNU General Public License
% along with this script.  If not, see <https://www.gnu.org/licenses/>.

\usepackage{amsmath,amssymb,gensymb,mathtools}
\usepackage[nottoc]{tocbibind}
%opening
\title{Formula for calculating $(P_j)_{1,1}$}
\author{Simon Crase}
\newcommand\numberthis{\addtocounter{equation}{1}\tag{\theequation}}
\begin{document}

\maketitle
\tableofcontents
\begin{abstract}
Check calculations from \cite{ssd}
\end{abstract}

\section{The Orbit in Space}

From  \cite[(2.119) and (2.120)]{ssd}:
\begin{align*}
P_1 =& \begin{pmatrix}
\cos \omega & - \sin \omega & 0\\
\sin \omega & cos \omega & 0\\
0 & 0  & 1
\end{pmatrix}
\text{, }
P_2 = \begin{pmatrix}
1 & 0 & 0\\
0 & \cos I & - \sin I\\
0 & \sin I & \cos I
\end{pmatrix}
P_3 =& \begin{pmatrix}
\cos \Omega & - \sin \Omega & 0\\
\sin \Omega & cos \Omega & 0\\
0 & 0  & 1
\end{pmatrix}
\end{align*}
$\varpi$
\begin{align*}
P_j \triangleq& P_3 P_2 P_1\\
 =& P_3 \begin{pmatrix}
1 & 0 & 0\\
0 & \cos I & - \sin I\\
0 & \sin I & cos I
\end{pmatrix}
\begin{pmatrix}
\cos \omega & - \sin \omega & 0\\
\sin \omega & cos \omega & 0\\
0 & 0  & 1
\end{pmatrix} \\
=& \begin{pmatrix}
\cos \Omega & - \sin \Omega & 0\\
\sin \Omega & cos \Omega & 0\\
0 & 0  & 1
\end{pmatrix}
\begin{pmatrix}
\cos \omega & -\sin \omega & 0\\
\cos I \sin \omega & \cos I \cos \omega & - \sin I\\
\sin I \sin \omega & \sin I \cos \omega & cos I
\end{pmatrix}\\
=&\begin{pmatrix*}[l]
\cos \Omega \cos \omega - \sin \Omega \cos I \sin \omega&-\cos \Omega \sin \omega - \sin \Omega \cos I \cos \omega& \sin \Omega \sin I\\
\sin \Omega \cos \omega + \cos \Omega \cos I \sin \omega& - \sin \Omega \sin \omega + \cos \Omega \cos I \cos \omega& - \cos \Omega \sin I \\
\sin I \sin \omega & \sin I \cos \omega & cos I
\end{pmatrix*}\numberthis \label{eq:Pj}
\end{align*}

\begin{align*}
\begin{bmatrix*}[c]
X\\
Y\\
Z
\end{bmatrix*}
=&P_3 P_2 P_1\begin{bmatrix*}[c]
r \cos f\\
r \sin f\\
0
\end{bmatrix*}\\
=&r \begin{bmatrix*}[l]
(\cos \Omega \cos \omega - \sin \Omega \cos I \sin \omega) \cos f -(\cos \Omega \sin \omega + \sin \Omega \cos I \cos \omega) \sin f\\
(\sin \Omega \cos \omega + \cos \Omega \cos I \sin \omega) \cos f - (\sin \Omega \sin \omega - \cos \Omega \cos I \cos \omega) \sin f\\
\sin I \sin \omega \cos f + \sin I \cos \omega \sin f
\end{bmatrix*}\\
=&r \begin{bmatrix*}[l]
\cos \Omega (\cos \omega \cos f-\sin \omega \sin f) -\sin\Omega \cos I (\sin \omega \cos f + \cos \omega \sin f)\\
\sin \Omega (\cos \omega \cos f - \sin \omega \sin f) + \cos\Omega \cos I (\sin \omega \cos f + \cos \omega \sin f) \\
\sin I (\sin \omega \cos f + \cos \omega \sin f)
\end{bmatrix*}\\
=&r \begin{bmatrix*}[l]
\cos \Omega \cos (\omega + f) - \sin \Omega \cos I \sin(\omega + f)\\
\sin \Omega \cos (\omega + f) + \cos \Omega \cos I \sin (\omega + f)\\
\sin I \sin (\omega + f)
\end{bmatrix*} \numberthis \label{eq:XYZ}
\end{align*}
Equation (\ref{eq:XYZ}) agrees with \cite[(2.122)]{ssd}.

The values for $\Omega$, $I$, and $\omega$ in the following table were taken from \cite[Page 51]{ssd}, and the remaining columns were calculated by Excel.
\begin{center}
	\begin{tabular}{|l|l|l|l|l|}
		\hline
		&&	radians&$\cos$&$\sin$\\ \hline
		$\Omega$&100.535\degree	&1.754666763&	-0.182836128&	0.983143403\\ \hline
		$I$&1.30537\degree&	0.022783004&	0.999740479&	0.022781034\\ \hline
		$\omega$&14.7392\degree&	0.257247569&	0.967093913&	0.25441966\\ \hline
	\end{tabular}
\end{center}

Substituting in (\ref{eq:Pj})

\begin{align*}
(P_j)_{1,1} =& \cos \Omega \cos \omega - \sin \Omega \cos I \sin \omega \\
=& -0.182836128 \times 	0.967093913 - 0.983143403 \times 0.999740479 \times 	0.25441966\\
=& -0.176819706 - 0.250066096\\
=& -0.426885803
\end{align*}

\section {3 Body}

From \cite[(24)]{kotovych2002exactly}
\begin{align*}
	V =& -G \sum_{i<j}\frac{m_i m_j}{r_{jk}}\\
	\partial_r V =& G \sum_{i<j}\frac{\partial_r r_{jk}}{r_{jk}^2}\\
	\partial_{\theta} V =& G \sum_{i<j}\frac{\partial_{\theta} r_{jk}}{r_{jk}^2}\\
	\partial_{\rho} V =& G \sum_{i<j}\frac{\partial_{\rho} r_{jk}}{r_{jk}^2}\\
	\partial_{\Theta} V =& G \sum_{i<j}\frac{\partial_{\Theta} r_{jk}}{r_{jk}^2}
\end{align*}

Using matrix notation, we define
\begin{align*}
	\nabla V \triangleq& \begin{pmatrix}
		\partial_r V\\
		\partial_{\theta} V\\
		\partial_{\rho} V\\
		\partial_{\Theta} V
	\end{pmatrix}\\
	S \triangleq& \begin{pmatrix}
		\partial_r r_{12}& \partial_r r_{23} & \partial_r r_{31}\\
		\partial_{\theta} r_{12}& \partial_{\theta} r_{23} & \partial_{\theta} r_{31}\\
		\partial_{\rho} r_{12}& \partial_{\rho} r_{23} & \partial_{\rho} r_{31}\\
		\partial_{\Theta} r_{12}& \partial_{\Theta} r_{23} & \partial_{\Theta} r_{31}
	\end{pmatrix}\\
	T \triangleq& \begin{pmatrix}
		\frac{m_1 m_2}{r_{12}^2}\\
		\frac{m_2 m_3}{r_{23}^2}\\
		\frac{m_1 m_3}{r_{13}^2} 		
	\end{pmatrix}\text{, then we can write}\\
	\nabla V =& G S T
\end{align*}
From  \cite[(27a)]{kotovych2002exactly}
\begin{align*}
	\vec{r_2} - \vec{r_1} =& \vec{r} \textit{, whence}\\
	r_{12} =& r \\
		S =& \begin{pmatrix}
		1& \partial_r r_{23} & \partial_r r_{31}\\
		0& \partial_{\theta} r_{23} & \partial_{\theta} r_{31}\\
		0& \partial_{\rho} r_{23} & \partial_{\rho} r_{31}\\
		0& \partial_{\Theta} r_{23} & \partial_{\Theta} r_{31}
	\end{pmatrix}\\
\end{align*}
From  \cite[(27b)]{kotovych2002exactly}
\begin{align*}
	\vec{r_3} - \vec{r_1} =& \vec{\rho} + \frac{m_2}{\mu}\vec{r} \numberthis \label{eq:r_3:r_1}\\
	=&\big[\rho \cos{\Theta}+\frac{m_2}{\mu}r \cos{\theta}\big] \vec{i_1} +\big[\rho \sin{\Theta}+\frac{m_2}{\mu}r \sin{\theta}\big] \vec{i_2}\\
	r_{13}^2 =& \big[\rho \cos{\Theta}+\frac{m_2}{\mu}r \cos{\theta}\big]^2 + \big[\rho \sin{\Theta}+\frac{m_2}{\mu}r \sin{\theta}\big]^2
\end{align*}

\begin{align*}	
	r_{13} \partial_r r_{13} =& \big[\rho \cos{\Theta}+\frac{m_2}{\mu}r \cos{\theta}\big]\frac{m_2}{\mu} \cos{\theta} + \big[\rho \sin{\Theta}+\frac{m_2}{\mu}r \sin{\theta}\big] \frac{m_2}{\mu} \sin{\theta}\\
	=& \frac{m_2}{\mu} \big[\cos{(\Theta-\theta)} \rho + \frac{m_2}{\mu} r\big]\\
	r_{13} \partial_{\theta} r_{13} =&-\big[\rho \cos{\Theta}+\frac{m_2}{\mu}r \cos{\theta}\big] \frac{m_2}{\mu}r \sin {\theta}+ \big[\rho \sin{\Theta}+\frac{m_2}{\mu}r \sin{\theta}\big]\frac{m_2}{\mu}r \cos{\theta}\\
	=& \frac{m_2}{\mu} \sin(\Theta-\theta) r \rho \\
	r_{13} \partial_{\rho} r_{13} =&\big[\rho \cos{\Theta}+\frac{m_2}{\mu}r \cos{\theta}\big] \cos{\Theta} + \big[\rho \sin{\Theta}+\frac{m_2}{\mu}r \sin{\theta}\big] \sin{\Theta}\\
	=& \rho + \big[\frac{m_2}{\mu}\big]^2 r \cos(\theta-\Theta)\\
	r_{13} \partial_{\Theta} r_{13} =&-\big[\rho \cos{\Theta}+\frac{m_2}{\mu}r \cos{\theta}\big] \rho \sin{\Theta}+ \big[\rho \sin{\Theta}+\frac{m_2}{\mu}r \sin{\theta}\big] \rho\cos{\Theta}\\
	=& \big[\frac{m_2}{\mu}\big]^2 r \rho \sin(\theta-\Theta)
\end{align*}
From  \cite[(27c)]{kotovych2002exactly}
\begin{align*}
	\vec{r_3} - \vec{r_2} =& \vec{\rho} - \frac{m_1}{\mu}\vec{r} 
\end{align*}
The RHS is identical to \eqref{eq:r_3:r_1} if we substitute $-m_1$ for $m_2$.
\bibliographystyle{unsrt}
\bibliography{astro}

\end{document}
